% alg-p1-04: 实数与数轴一一对应——用勾股定理在数轴上标出 sqrt(2)
\documentclass[tikz,border=4pt]{standalone}
\usepackage{ctex}
\usepackage{amsmath}
\usetikzlibrary{arrows.meta, calc}
\begin{document}
\begin{tikzpicture}[scale=1.1, thick]

  % 数轴
  \draw[-{Stealth[length=5pt]}] (-0.4,0) -- (3.8,0);

  % 刻度
  \foreach \x in {0,1,2,3} {
    \draw (\x, 3pt) -- (\x, -3pt) node[below, font=\small] {$\x$};
  }
  \node[below left, font=\small] at (0,0) {$O$};

  % ---- 构造正方形：以 (0,0)-(1,0) 为底边 ----
  \draw[gray, dashed] (0,0) -- (1,0) -- (1,1) -- (0,1) -- cycle;
  % 边长标注
  \node[left, font=\footnotesize, gray] at (0,0.5) {$1$};
  \node[below, font=\footnotesize, gray] at (0.5,0) {$1$};

  % 对角线（长度 sqrt(2)）
  \draw[red, thick] (0,0) -- (1,1);
  \node[above left, font=\footnotesize, red] at (0.5,0.5) {$\sqrt{2}$};

  % 以原点为圆心，半径 sqrt(2) 作弧，交正半轴
  \draw[blue, thick, -] (1,1) arc ({atan2(1,1)}:0:{sqrt(2)});

  % 在数轴上标出 sqrt(2) 的位置
  \filldraw[blue] ({sqrt(2)},0) circle (2.5pt);
  \draw[blue] ({sqrt(2)}, 3pt) -- ({sqrt(2)}, -3pt);
  \node[below, font=\small, blue] at ({sqrt(2)}, -0.05) {$\sqrt{2}$};

  % 小圆弧箭头提示（从斜边端点到数轴点）
  \draw[-{Stealth[length=4pt]}, blue, dashed, thin]
    (1,1) arc ({atan2(1,1)+10}:{5}:{sqrt(2)});

  % 说明文字
  \node[font=\footnotesize, align=left] at (2.5, 0.85)
    {每个实数\\对应数轴\\唯一一点};

\end{tikzpicture}
\end{document}
